\section{The shape of the Lamperti density}\label{sec:shape}

We prove Theorem~\ref{thm:shape}. Let $\alpha_c$ be the unique root of
$\cos(\pi\alpha/2)=\alpha$ in $(0,1)$. The function $\alpha\mapsto\cos(\pi\alpha/2)-\alpha$
decreases strictly from $1$ to $-1$ on $[0,1]$, so $\alpha_c$ exists, and
$\cos(\pi\alpha/2)\ge\alpha$ exactly when $\alpha\le\alpha_c$. Numerically
$\alpha_c=0.594611644056\ldots$.

\begin{proof}[Proof of Theorem~\ref{thm:shape}]
Let $b=1$ and $c=\cos\pi\alpha$. By~\eqref{eq:L-cosh} with $\beta=0$,
$L_{\alpha,1}(x)=\frac{\sin\pi\alpha}\pi\,\phi(t)$ with
\[
 \phi(t)=\frac{1+\cosh t}{\cosh\alpha t+c},\qquad t=\log\frac{1-x}x .
\]
The map $x\mapsto t$ is a decreasing bijection from $(0,1)$ onto $\R$, so the critical
points of $L_{\alpha,1}$ correspond to those of $\phi$, with the same type. The function
$\phi$ is even, so $t=0$, that is $x=1/2$, is always a critical point. Differentiating,
\[
 \phi'(t)=\frac{(1+\cosh t)\bigl[\tanh(t/2)(\cosh\alpha t+c)-\alpha\sinh\alpha t\bigr]}{(\cosh\alpha t+c)^2},
\]
since $\sinh t=(1+\cosh t)\tanh(t/2)$. For $t>0$ the sign of $\phi'(t)$ is therefore the
sign of
\[
 \chi(t)=\cosh\alpha t+c-\alpha\sinh(\alpha t)\coth(t/2)=G(t)+c,\qquad
 G(t)=\cosh\alpha t-\alpha\sinh(\alpha t)\coth(t/2).
\]
We claim that $G$ is strictly increasing on $(0,\infty)$. Put $s=t/2$. Using
$1+1/(2\sinh^2s)=\cosh t/(2\sinh^2s)$ and $\coth s=\sinh t/(2\sinh^2s)$ we get
\[
 G'(t)=\alpha\sinh\alpha t-\alpha^2\cosh(\alpha t)\coth s+\frac{\alpha\sinh\alpha t}{2\sinh^2s}
 =\frac{\alpha\,n(t)}{2\sinh^2(t/2)},
\]
with $n(t)=\sinh(\alpha t)\cosh t-\alpha\cosh(\alpha t)\sinh t$. Now $n(0)=0$ and
\[
 n'(t)=(1-\alpha^2)\sinh(\alpha t)\sinh t>0\qquad(t>0),
\]
so $n>0$ and $G'>0$ on $(0,\infty)$. This proves the claim. As $t\downarrow0$,
$\alpha\sinh(\alpha t)\coth(t/2)\to2\alpha^2$, so $G(0+)=1-2\alpha^2$. As $t\to\infty$,
$G(t)=\frac{1-\alpha}2e^{\alpha t}(1+o(1))\to\infty$. Hence $\chi$ increases strictly
from
\[
 \chi(0+)=1+\cos\pi\alpha-2\alpha^2=2\bigl(\cos(\pi\alpha/2)-\alpha\bigr)\bigl(\cos(\pi\alpha/2)+\alpha\bigr)
\]
to $+\infty$. The second factor is positive.

If $\cos(\pi\alpha/2)\ge\alpha$, then $\chi>0$ on $(0,\infty)$. So $\phi$ increases strictly
on $[0,\infty)$ and, being even, decreases strictly on $(-\infty,0]$. Its only critical
point is the minimum at $0$, and $L_{\alpha,1}$ decreases strictly on $(0,1/2]$ and
increases strictly on $[1/2,1)$.

If $\cos(\pi\alpha/2)<\alpha$, then $\chi$ has exactly one zero $t^*>0$, with $\chi<0$
on $(0,t^*)$ and $\chi>0$ on $(t^*,\infty)$. So $\phi$ has a strict local maximum at
$0$ and strict minima at $\pm t^*$, and no other critical point. In the variable $x$
these are the maximum at $1/2$ and the minima at $x=1/(1+e^{\pm t^*})$.
\end{proof}

Near the threshold the two minima split off from $1/2$ like a pitchfork.

\begin{proposition}[the minima near the threshold]\label{prop:pitchfork}
For $\alpha>\alpha_c$ let $t^*_\alpha>0$ be the zero of $\chi$ in the proof of
Theorem~\ref{thm:shape}, so the minima of $L_{\alpha,1}$ are at $1/(1+e^{\pm t^*_\alpha})$. Then
\[
 t^*_\alpha\sim\Bigl(\frac{3\,(2\alpha^2-1-\cos\pi\alpha)}{\alpha^2(1-\alpha^2)}\Bigr)^{1/2}
 \sim\Bigl(\frac{12\,(1+\frac\pi2\sin\frac{\pi\alpha_c}2)}{\alpha_c(1-\alpha_c^2)}\,(\alpha-\alpha_c)\Bigr)^{1/2}
 \qquad(\alpha\downarrow\alpha_c).
\]
In particular the minima are at distance $t^*_\alpha/4+O((t^*_\alpha)^3)$ from $1/2$.
\end{proposition}

\begin{proof}
Expanding at $t=0$, $n(t)=\frac13\alpha(1-\alpha^2)t^3+O(t^5)$, so
$G'(t)=\frac23\alpha^2(1-\alpha^2)\,t+O(t^3)$ and
\[
 \chi(t)=\chi(0+)+\kappa_\alpha t^2\bigl(1+r(t,\alpha)\bigr),\qquad
 \kappa_\alpha=\tfrac13\alpha^2(1-\alpha^2),
\]
with $|r(t,\alpha)|\le Ct^2$ for $|t|\le1$ and $\alpha$ in a neighbourhood of $\alpha_c$, since
$\chi$ is an analytic function of $(t,\alpha)$ near $(0,\alpha_c)$ and even in $t$. As
$\alpha\downarrow\alpha_c$ we have $\chi(0+)\uparrow0$. The difference $\chi(1)-\chi(0+)=G(1)-G(0+)$
is positive and continuous in $\alpha$, so $\chi(1)>0$ for $\alpha$ close to $\alpha_c$, and the zero
$t^*_\alpha$ lies in $(0,1)$. The same argument shows that $t^*_\alpha\to0$. For each fixed
$\varepsilon\in(0,1)$ the difference $\chi(\varepsilon)-\chi(0+)=G(\varepsilon)-G(0+)$ is positive and
continuous in $\alpha$, so $\chi(\varepsilon)>0$ and $t^*_\alpha<\varepsilon$ for $\alpha$ close to
$\alpha_c$. Hence $r(t^*_\alpha,\alpha)=O((t^*_\alpha)^2)\to0$. The equation $\chi(t^*_\alpha)=0$ gives
$(t^*_\alpha)^2(1+O((t^*_\alpha)^2))=-\chi(0+)/\kappa_\alpha$, which is the first equivalence. For the
second, $-\chi(0+)=2(\alpha-\cos\frac{\pi\alpha}2)(\alpha+\cos\frac{\pi\alpha}2)$, and as
$\alpha\downarrow\alpha_c$ the first factor is
$(1+\frac\pi2\sin\frac{\pi\alpha_c}2)(\alpha-\alpha_c)(1+o(1))$ and the second tends to $2\alpha_c$.
Finally $1/(1+e^{t})=\frac12-\frac t4+O(t^3)$.
\end{proof}

The first expression in Proposition~\ref{prop:pitchfork} is within $0.01\%$ of $t^*_\alpha$ at
$\alpha=0.5947$, where $t^*_\alpha=0.0790$, and within $1.5\%$ at $\alpha=0.61$, where
$t^*_\alpha=1.018$ and the minima sit at $x=0.2655$ and $0.7345$.

At $\alpha=1/2$ the theorem gives the arcsine density
$L_{1/2,1}(x)=1/(\pi\sqrt{x(1-x)})$, with its minimum at $1/2$. As $\alpha\to1$ the two
minima move to the ends, since $\chi(t)\to\cosh t-1-\sinh(t)\coth(t/2)=-2$ for every
fixed $t$ while $\chi(t)\to\infty$ as $t\to\infty$ for each $\alpha<1$.

Margolin and Barkai~\cite[Section~III]{margolinbarkai2006} state that the transition
between the U shape and the W shape of $L_{\alpha,1}$ happens at $\alpha=0.5946\ldots$,
in agreement with Theorem~\ref{thm:shape}. In the arXiv version of their paper we found
the value but not its derivation. Zaburdaev, Denisov and Klafter~\cite[Section~II.A.3]{zdk2015}
note that the propagator of the ballistic L\'evy walk is U-shaped and becomes W-like
for $\alpha\gtrsim0.6$.

\begin{remark}[unequal tails]\label{rem:unequal}
For $b\ne1$ the same computation with~\eqref{eq:L-cosh} shows that the sign of the
derivative in $t$ is the sign of $\tanh(t/2)(\cosh(\alpha t-\beta)+c)-\alpha\sinh(\alpha t-\beta)$,
which tends to $-\infty$ as $t\to-\infty$ and to $+\infty$ as $t\to\infty$. So the number
of critical points of $\La$ is odd when they are all simple. Numerically it is $1$ or
$3$. For instance at $\alpha=0.7$ it is $3$ for $b=1$ and $b=2$ and $1$ for $b=3$. We do
not have a proof for $b\ne1$.
\end{remark}
